\chapter*{Цели и задачи}
\addcontentsline{toc}{chapter}{Введение}
\label{ch:intro_lab7}

\textbf{Цель лабораторной работы:} научиться применять разработанный пайплайн для тиражирования кода с целью решения широкого круга задач машинного обучения.

\textbf{Основные задачи:}
\begin{itemize}
    \item получение навыков рефакторинга кода в проектах машинного обучения;
    \item получение навыков определения ключевых признаков в задачах машинного обучения;
    \item получение навыков реализации ключевых стратегий оптимизации моделей регрессии.
\end{itemize}

В качестве набора данных для выполнения индивидуального задания используется набор данных \textbf{Abalone} (морское ушко), который содержит информацию о возрасте моллюсков, определяемом по количеству колец на раковине. Данный набор содержит как числовые, так и категориальные признаки.

\endinput